\chapter*{Terms used in this dissertation}
\addcontentsline{toc}{chapter}{Terms used in this dissertation}

The first three are measured quantities this dissertation turns on, and they have
different denominators; confusing them is the easiest way to misread the results. The
one-line forms below are for reference while reading. Their formulae, denominators and
caveats are given where they are derived, in \S\ref{sec:des-shield}.

\begin{description}[leftmargin=0pt, itemsep=2pt]

\item[Authored share.] Of logged decision ticks, the fraction whose served phase
originated with the model. The accountability number.

\item[Proposal acceptance.] Of ticks where the model returned a well-formed proposal,
the fraction served. Not a measure of authorship: the model is not consulted on every
tick.

\item[Divergence rate.] Of the decisions the model authored, the fraction that differ
from MaxPressure's choice.

\item[Congestion gate.] The rule deciding whether the model is consulted at all. Below a
threshold of halting demand, MaxPressure serves the tick directly.

\item[Anti-starvation floor.] A best-effort minimum-service rule that force-serves any
approach passed over beyond a bounded window, overriding both the model and
MaxPressure. It is suspended during an admissible emergency preemption, and records a
violation when two phases reach the threshold together. It authors a substantial
minority of logged decision ticks.

\item[Shield.] The safety layer: phase masking, MaxPressure substitution for invalid
output, and the floor. A sound heuristic composite, not correct-by-construction in the
formal sense (\S\ref{sec:bg-shields}).

\item[Substitution path.] The part of the shield replacing an invalid or absent
proposal. Measured at $0.0\%$ of decisions almost everywhere, which is what stops
``the model matches MaxPressure'' from being circular.

\item[MaxPressure.] The classical adaptive baseline. It serves the phase of greatest
\emph{pressure}, upstream queue minus the downstream queue it would discharge into, so
it declines to serve a long queue whose exit is blocked.

\item[Fixed-time floor.] The naive comparator: each network's own default cyclic
program, with no adaptive controller. The standard every deployed junction already meets.

\item[Longest-queue-first.] A third classical baseline serving the largest upstream
queue with no downstream term. A different rule from MaxPressure, though measured
indistinguishable from it here.

\item[Teleport.] SUMO's gridlock-escape mechanism, which removes a vehicle stuck beyond
a timeout. A high count means the network is failing rather than flowing.

\item[Demand calibration.] Subsampling demand until the network clears, so controller
differences are measurable rather than compressed by gridlock.

\item[QLoRA.] The fine-tuning method: low-rank adapters trained against a 4-bit
quantised base, which makes training feasible on the GPU that also serves the model.

\item[int4 ONNX.] The served weight format, compiled to fit the 6\,GB envelope.

\item[Foundry Local.] Microsoft's on-device serving runtime, and a fixed constraint of
this project rather than a choice.

\item[Ed25519 hash chain.] The audit structure. Each decision record is signed and
linked to its predecessor by a hash, so later alteration breaks the chain detectably.

\item[LOTO.] Leave-one-topology-out: a control student trained with one network's states
removed, testing whether gains are topology memorisation.

\item[Stock, ft.] \emph{Stock} is the same base model through the identical compilation
pipeline but not fine-tuned; \emph{ft} is the fine-tuned student. Comparisons are always
against the identically-compiled control, never the vendor artifact.

\end{description}
